\subsection{Forecasting}\label{sec:forecasting}
\gretlpy\footnote{This section is based on \cite{gretl_manual}.}
offers predictive capabilities in the same way as
\gretl. In the simplest case, the ``predictions''
of interest are simply the (within-sample) fitted
values from an econometric model. If the model in
question takes the form of a system of equations,
\gretlcmd{\$yhat} returns a matrix, each column of
which contains the fitted values for a particular
dependent variable. To generate forecasts, the user
should use the \pythcode{fcast()} method of a
\classpymodel instance (see \ref{met:gretlmodel::fcast}).

By default, the forecast range coincides with the
currently defined sample range. If this range remains
unchanged following estimation of the model in
question, the forecast will be ``within sample'' and
will essentially duplicate the information available
via retrieval of the fitted values.

A common situation is that a model is estimated over
a given sample and forecasts are then required for a
subsequent out-of-sample range. The simplest way to
accomplish this is by using the \pythcode{out\_of\_sample}
option of the \pythcode{.fcast()} method. For example,
suppose we have a quarterly time-series dataset
containing observations from 1980:1 to 2008:4, of
which the last four are to be reserved for forecasting:
\begin{lstlisting}[style=pythonstylefoot]
import gretl

d = gretl.get_data('denmark.gdt', quiet=True)
d.new_list('Xes', 'LRY IBO IDE')

# reserve the last 4 observations
d.sample('1974:1', '1986:3', quiet=True)

m = gretl.ols('LRM const Xes').fit()
m.fcast(out_of_sample=True)
\end{lstlisting}\index{get\_data()}\index{new\_list()}\index{sample()}\index{ols()}\index{fcast()}
This code generates a forecast from 1986:4 to 1987:3.
Figure~\ref{fig:OLS_based_forecasts} shows results of above script.

\begin{figure}[h]
	\centering
	\includegraphics[width=0.6\columnwidth,keepaspectratio]{forecasting_ols}
	\caption{OLS-based forecasts}
	\label{fig:OLS_based_forecasts}
\end{figure}

Alternatively, the forecast range can be controlled
in the following ways:
\begin{itemize}
	\item Use the \pythcode{.smpl()} method to adjust the
	sample range prior to invoking \pythcode{.fcast()}.
	\item Use the optional \textit{startobs} and
	\textit{endobs} arguments of \pythcode{fcast()}.
	These values set the forecast range independently
	of the sample range.
\end{itemize}

The interface for forecasting from a VAR is similar to
that for a single equation. Here's an example code:
\begin{lstlisting}[style=pythonstylefoot]
import gretl

# open sample data file
d = gretl.get_data('sw_ch14.gdt', quiet=True)

# generate the "inflation" series
d.new_series('INFL = 100 * sdiff(log(PUNEW))')

# put last year aside for out-of-sample forecast
d.sample(1, -4, quiet=True)

# estimate a 5-lag VAR
m = gretl.var(5, 'LHUR INFL', silent=True).fit()

# retrieve fitted values (note: result is a 2-column matrix)
YH = m.get_accessor('$yhat')

# perform out-of-sample forecast (both versions)
m.fcast('LHUR', static=True, out_of_sample=True)

# note that omission of the variable specification means "all"
m.fcast(dynamic=True, out_of_sample=True)
\end{lstlisting}\index{get\_data()}\index{new\_series()}\index{sample()}\index{var()}\index{get\_accessor()}\index{fcast()}
Figure~\ref{fig:VAR_based_forecasts} shows results of above script.

\begin{figure}
	\centering
	\includegraphics[height=\textheight,keepaspectratio]{forecasting_var}
	\caption{VAR-based forecasts}
	\label{fig:VAR_based_forecasts}
\end{figure}

Using a cointegrated VAR model with \gretl's \gretlcmd{vecm}
does not change the way a forecast is obtained afterwards.
The VECM can be internally represented as a VAR (in levels)
that automatically contains the reduced-rank restrictions of
cointegration, and this VAR form is then used to calculate
the forecasts. Providing forecast standard errors and the
associated confidence bands is also straightforward since only
the innovation uncertainty is captured in those. This ease
of use also carries over to the situation when a VECM with
additional exogenous terms is used for forecasting---provided
that future values of the exogenous variables are specified, of course.
